% =====================================================================
%  translator_note.tex —— 译者说明（放在目录之后、摘要之前）
% =====================================================================
\section*{译者说明}
\addcontentsline{toc}{section}{译者说明}

\textbf{底本。}本书译自 Dan Xie, \emph{General Argyres-Douglas Theory}, arXiv:1204.2270v2 [hep-th]（2016 年 7 月 12 日；Prepared for submission to JHEP），共 54 个印刷页，含插图 20 幅、表格 4 张、参考文献 67 条。文中「p.$N$」一律指底本的印刷页码，与原 PDF 的页序相差 2。

\textbf{版式。}公式编号、图号、表号、文献编号全部照底本原样，不重排、不跳号。插图按底本重绘为黑白线稿，图内标注保留底本写法。行文引用公式时写「式 (n.m)」，不另造编号。

\textbf{术语。}底本使用的术语中，有相当一部分在中文学界没有通行译名，或者通行译名与英文原词的实际所指不符。对这些术语，我们只采用两种处置：一是保留英文原文（例如 abelianization、spray、class S、magical、lamination space），二是在仔细讨论后定出一个中译，并在首次出现处用译注交代其英文与所指。凡属后一种情形，其定案都不是临时决定的。这一领域的核心译名（非正则奇点、正则奇点、穿刺点、极点的阶、牛顿多边形、Stokes 扇区与射线、单值性、Hitchin 纤维化等）沿用本库同一主题既有译本已冻结的译法；尚无定名者，先由两位术语译手各自出提案、逐条互审，再由译者定案并记入工程内的术语表。全部术语的最终写法以该表为准，正文不得另立译名。

\textbf{原文笔误。}底本有若干处笔误与前后不一致（公式中的幂次、分母、下标，图注与正文对不上，符号重号等）。我们的原则是：正文照录原文，不加改动；凡影响理解者，用【原文如此】标出原样，再在译注里给出正确形式与判定依据（通常是由相邻公式反算，或由底本自身的另一处写法互证）。凡属不影响理解的排印瑕疵，径行订正，不出注。

\textbf{译注。}全书有三类标注。

\begin{itemize}
  \item 【译注】：解释性补充。每条都带来源，写作 \verb|[来源：…]| 或 \verb|[依据：…]|，其中来源按 A（底本、作者其他论文、经典教材）、B（nLab、Stacks Project、ProofWiki 等）、I（本库既往译本）三级标注；确实没有外部依据的，写作 \verb|[译者按（无外部依据）：…]|。
  \item 【路线】：交代作者的论证路线——为什么先证这一条、哪条路被放弃、某处结论依赖什么前提。这类注带底本页码。
  \item 【原文如此】：指出底本笔误之处，正文未作改动。
\end{itemize}

\textbf{阅读建议。}第 3 节是技术核心：它给出 Hitchin 方程非正则奇点的分类，后面各节都在用这套分类。只关心物理结果的读者可以跳过第 3.1 节的格点与杨图细节，先读第 4 至 7 节；但第 6.1 节把非正则奇点压成三类的那一步论证（式 (6.6) 处的矛盾）是全篇的枢纽，不宜略过。
